\begin{中文摘要}{端云协同；大语言模型；投机解码；投机采样；Android 推理}
随着大语言模型在智能问答、文本生成和辅助编程等场景中的应用不断扩大，模型推理对算力、存储和时延的需求持续增加。移动终端交互便捷、部署灵活，但受限于处理器性能、内存容量和功耗预算，难以独立承担高质量大模型推理任务；远端设备计算能力较强，但纯远程推理会受到网络时延和调用开销影响。投机解码通过草稿模型和目标模型协同生成，使目标模型在一次验证中推进多个词元，为端云协同推理提供了可利用的任务划分方式。

本文面向“移动端草稿生成、桌面端目标验证”的协同推理场景，设计并实现了一套端云协同投机解码系统。系统以 Android 端作为草稿生成侧，以桌面端推理服务作为目标验证侧，在远程推理通信基础上引入投机解码协议、会话管理机制以及接受、拒绝和修正逻辑。针对跨设备协同中的状态同步、词元空间一致性和验证开销控制问题，系统设计了草稿会话与目标会话两类状态对象，支持统一真实词元空间下的线性提案、树形候选和目标验证，并通过运行时连续性维护和请求字段精简降低单轮协同开销。

实验部分以云侧本地单进程投机解码和云侧本地双进程投机解码作为对照，对 Android--桌面端跨设备投机解码方案进行评估。实验结果表明，所实现系统能够完成稳定的多轮投机会话，跨设备方案在长输出任务中取得了接近云侧本地双进程方案的吞吐表现，验证了移动端草稿生成与桌面端目标验证协同工作的可行性。进一步分析显示，跨设备方案的主要性能差距来自 Android 草稿侧状态维护、跨设备协同通信以及少量接受率差异。
\end{中文摘要}

\begin{英文摘要}{collaborative inference; large language model; speculative decoding; speculative sampling; Android inference}
With the adoption of large language models in question answering, text generation, and coding assistance, inference places high demands on computation, memory, and latency. Mobile devices provide convenient interaction and flexible deployment, but their limited compute capability, memory capacity, and power budget make it difficult to support high-quality large-model inference independently. Remote devices provide higher computing capacity, while pure remote inference is affected by network latency and invocation overhead. Speculative decoding coordinates a draft model and a target model so that multiple tokens can be advanced through one target-side verification step.

This work designs and implements a collaborative speculative decoding system for an Android-desktop inference setting. In the system, the Android device acts as the draft generation side, while the desktop inference service acts as the target verification side. On top of remote inference communication, the system introduces a speculative decoding protocol, session management, and acceptance, rejection, and correction logic. To address state synchronization, token-space consistency, and verification overhead in cross-device execution, the system maintains separate draft and target sessions, supports linear proposals, tree-based candidates, and target verification in a unified real-token space, and reduces per-round coordination overhead through runtime continuity and compact request fields.

The experiments use cloud-side local single-process speculative decoding and cloud-side local dual-process speculative decoding as baselines to evaluate the Android--desktop cross-device scheme. The results show that the implemented system can support stable multi-round speculative sessions, and the cross-device scheme achieves throughput close to the cloud-side local dual-process setting on long-output generation. This verifies the feasibility of coordinating mobile-side draft generation with desktop-side target verification. Further analysis shows that the remaining performance gap mainly comes from Android-side state maintenance, cross-device coordination overhead, and a small acceptance-rate difference.
\end{英文摘要}
